\Needspace{18\baselineskip}
\section{Finite carriers and concrete interfaced graphs}

We now turn to actual graph objects. The constructions in the following sections
prove their own incidence and typing obligations. They do not assume the specimen
laws used in the preceding abstract argument. We work with small finite vertex
and edge carriers; operation labels may live in a larger universe.

\Needspace{24\baselineskip}
\subsection{Enumerations and bijections}
\begin{definition}[Finite enumeration]\label{term:enumeration}
An enumeration of a type $A$ consists of a finite list $l$ and a proof that every
$a\in A$ belongs to $l$. Repetitions are allowed. This witnesses finiteness without
requiring decidable equality or a canonical choice of names.
\end{definition}
\lean{Graph/Finite}{9}{11}{Enumeration}
The empty type is enumerated by the empty list: completeness follows because
there is no element to check. The finite ordinal $\{0,\ldots,n-1\}$ is enumerated
in increasing order by Lean's finite range; its library membership theorem proves
completeness, including when $n=0$.
\lean{Graph/Finite}{15}{15}{Enumeration.empty}
\lean{Graph/Finite}{17}{17}{Enumeration.fin}
For a disjoint union $A+B$, concatenate the tagged enumerations of $A$ and $B$.
An element has exactly one tag, and completeness of the corresponding list places
it in the appropriate half of the concatenation.
\lean{Graph/Finite}{19}{25}{Enumeration.sum}
For a subset $\{a\in A\mid p(a)\}$, attach membership certificates to the enumeration
of $A$ and retain just entries satisfying $p$. Classical decidability of $p$
allows this filtering. Every subset element is retained because its certificate
proves the test true. The resulting definition is noncomputable for an arbitrary
proposition $p$; no algorithm deciding $p$ is claimed.
\lean{Graph/Finite}{27}{37}{Enumeration.subtype}
For a dependent sum $\sum_{a\in A}B_a$, enumerate each fibre, tag it with $a$, and
concatenate over the enumeration of $A$. A pair $(a,b)$ appears because $a$ occurs
in the outer list and $b$ in the enumeration of its fibre.
\lean{Graph/Finite}{39}{43}{Enumeration.sigma}

\begin{definition}[Bijection]\label{term:bijection}
A bijection $f:A\cong B$ consists of functions $f:A\to B$ and $f^{-1}:B\to A$
with both inverse equations, pointwise.
\end{definition}
\lean{Graph/Finite}{47}{51}{Bijection}
Identity uses the identity function twice. Inversion interchanges the two functions
and their inverse laws. Composition uses $g\circ f$ forwards and
$f^{-1}\circ g^{-1}$ backwards; cancelling the inner inverse pair and then the
outer pair proves each law.
\lean{Graph/Finite}{57}{57}{Bijection.refl}
\lean{Graph/Finite}{59}{60}{Bijection.symm}
\lean{Graph/Finite}{62}{66}{Bijection.trans}
Bijections extend to disjoint unions by applying the appropriate map inside each
tag. Case analysis on the tag reduces both inverse laws to the component laws.
\lean{Graph/Finite}{68}{73}{Bijection.sum}
Every bijection is injective: apply its inverse to an equality of images and
cancel. An enumeration transports along a bijection by mapping its forward
function over the list. For $b\in B$, the listed preimage $f^{-1}(b)$ maps to $b$,
proving completeness.
\lean{Graph/Finite}{75}{77}{Bijection.injective}
\lean{Graph/Finite}{79}{81}{Bijection.enumeration}
\begin{lemma}[Distinctness under restricted injectivity]
If a list $l$ has no repeated elements and $f$ is injective on elements of $l$,
then $\operatorname{map}(f,l)$ has no repeated elements.
\end{lemma}
\begin{proof}
Distinct positions in $l$ contain distinct elements. If their images were equal,
restricted injectivity would identify those elements, a contradiction. Lean
expresses this as preservation of the pairwise inequality predicate under mapping.
\end{proof}
\lean{Graph/Finite}{86}{92}{nodup_map_on}
This restricted version will be essential: replacement can identify an input
with an output, while remaining injective on each incidence list separately.

\Needspace{24\baselineskip}
\subsection{Ports, signatures, and graph invariants}
\begin{definition}[Signature and ports]\label{term:ports}
A signature $\sigma$ assigns an ordered input-label list and output-label list
to every operation. With internal carrier $V$ and boundary lengths $n,m$, vertices
are the disjoint union
\[
 \operatorname{Port}(V,n,m)=V+\operatorname{Fin}(n)+\operatorname{Fin}(m).
\]
The three constructors are internal, input, and output. In particular the two
boundary sets are disjoint even if their wire labels agree.
\end{definition}
\lean{Graph/Basic}{13}{15}{Signature}
\lean{Graph/Basic}{17}{21}{Port}
The predicates $\mathsf{IsInput}$ and $\mathsf{IsOutput}$ test the corresponding
constructor. A carrier map $f:V\to W$ acts only on internal vertices and fixes
each boundary position. It is injective whenever $f$ is: unlike constructors
cannot have equal images, and equality within the internal constructor reduces
to injectivity of $f$. The input and output predicates are preserved and reflected
by this map, by the same three-case analysis.
\lean{Graph/Basic}{27}{29}{Port.IsInput}
\lean{Graph/Basic}{31}{33}{Port.IsOutput}
\lean{Graph/Basic}{35}{38}{Port.map}
\lean{Graph/Basic}{40}{43}{Port.map_injective}
\lean{Graph/Basic}{45}{46}{Port.map_isInput}
\lean{Graph/Basic}{48}{49}{Port.map_isOutput}
\termdef{vertex-label}
Given boundary-label lists $I,O$ and internal labelling $\lambda:V\to L$, the label
of an internal vertex is $\lambda(v)$, of input $i$ is $I_i$, and of output $j$
is $O_j$. The finite-index types ensure that these lookups are defined.
\lean{Graph/Basic}{51}{54}{Port.label}

\begin{definition}[Concrete typed net]\label{term:net}
A net $G:\!I\to O$ has finite internal and edge carriers, labels for both, and
ordered source and target vertex lists $s(e),t(e)$ for each edge. It satisfies:
\begin{enumerate}
\item Each individual source list and target list has no repetitions.
\item A vertex belongs to the sources of at most one edge, and to the targets
of at most one edge.
\item A vertex belongs to some source list exactly when it is not an output;
it belongs to some target list exactly when it is not an input.
\item Mapping vertex labels over $s(e)$ and $t(e)$ gives precisely the input
and output lists of the signature of the edge label.
\end{enumerate}
\end{definition}
\lean{Graph/Basic}{61}{77}{Net}
These fields are part of a net's certificate of well-formedness. Coverage excludes
isolated vertices. It does not exclude cycles or a vertex belonging to both the
source and target of the same edge. The full vertex type and its label function
are abbreviations for the port type and the labelling just defined.
\lean{Graph/Basic}{84}{84}{Net.Vertex}
\lean{Graph/Basic}{86}{86}{Net.vertexLabel}
From forward source coverage, membership in $s(e)$ implies that a vertex is not
an output. From target coverage, membership in $t(e)$ implies that it is not an
input. Apply the appropriate equivalence to the witness edge $e$.
\lean{Graph/Basic}{88}{89}{Net.source_not_output}
\lean{Graph/Basic}{91}{92}{Net.target_not_input}
Taking lengths in the two typing equations shows that incidence-list lengths are
exactly the corresponding signature arities: mapping labels does not change length.
\lean{Graph/Basic}{94}{96}{Net.source_length}
\lean{Graph/Basic}{98}{100}{Net.target_length}
\termdef{indexed-port}
The $i$th source port of $e$ is its source list indexed by a position in the
signature's input list; the preceding length equality justifies the lookup.
Target ports are defined in the same way with the output list.
\lean{Graph/Basic}{102}{103}{Net.sourcePort}
\lean{Graph/Basic}{105}{106}{Net.targetPort}
Each indexed port belongs to its incidence list, by the list lookup membership
lemma. Within either list, equality of indexed ports implies equality of indices,
because that list has no duplicates; equality of finite ordinals follows from
equality of their underlying natural numbers.
\lean{Graph/Basic}{108}{108}{Net.sourcePort_mem}
\lean{Graph/Basic}{110}{110}{Net.targetPort_mem}
\lean{Graph/Basic}{112}{114}{Net.sourcePort_injective}
\lean{Graph/Basic}{116}{118}{Net.targetPort_injective}
The label of source port $i$ is the $i$th signature input label, and the label of
target port $i$ is the $i$th signature output label. Apply indexed lookup to the
appropriate list typing equation, then commute lookup with mapping.
\lean{Graph/Basic}{120}{123}{Net.sourcePort_label}
\lean{Graph/Basic}{125}{128}{Net.targetPort_label}
Finally, membership in a source list is equivalent to being its indexed port for
some signature input position. Given membership, choose its list index and
transport the bound along the arity equality. Conversely, indexed ports belong
to their lists. The target statement uses the output arity and the identical
membership argument.
\lean{Graph/Basic}{130}{137}{Net.mem_source_iff_port}
\lean{Graph/Basic}{139}{146}{Net.mem_target_iff_port}

\Needspace{24\baselineskip}
\subsection{Elementary graph constructors}
The ordered finite range has no duplicates: equality of its entries at two
positions gives equality of their underlying natural numbers and hence of the
positions.
\lean{Graph/Constructors}{5}{9}{finRange_nodup}
\begin{definition}[Corolla and empty graph]\label{term:corolla}
The corolla of operation $o$ has one edge labelled $o$, no internal vertices,
and all of its prescribed inputs and outputs in signature order. Its sources
are precisely the input constructors and its targets the output constructors.
\end{definition}
\begin{proof}[Well-formedness]
Finite ranges and their constructor maps give distinct incidence lists. There
is only one edge, so edge uniqueness is immediate. Every input occurs in its
source range and every output in its target range, and the opposite boundary
constructors cannot occur there. There are no internal vertices to check.
Mapping labels over either ordered range recovers the corresponding signature
list, proving typing.
\end{proof}
\lean{Graph/Constructors}{18}{52}{Net.corolla}
\termdef{empty-graph}
The empty graph has empty internal and edge carriers and empty boundaries. Each
incidence or label function has empty domain; all obligations follow by eliminating
the purported element of an empty type. A bare identity wire is not provided:
a single vertex serving as both boundary input and boundary output would violate
the disjoint-boundary representation used here.
\lean{Graph/Constructors}{56}{78}{Net.empty}
